\section{Introduction}
Consider a software team that uses AI agents to draft pull requests (PRs). Faster drafting can lengthen the review queue. The team may respond by routing routine changes to automated checks and reserving human review for difficult cases. To evaluate this change, it must decide which PRs to audit, how to combine new observations with earlier evidence, and who may revise the routing or audit rule. An apparent gain can reflect a better workflow, a different sample of reviewed cases, or more effective reuse of evidence. How can the team distinguish these explanations and retain changes that improve delivery?

We propose a modeling specification for this problem and evaluate it with an executable checker, a mapping of public PR records, and a controlled mechanism study. The running PR example connects the specification's fields to concrete organizational decisions; the simulation isolates evaluation and memory using prescribed task distributions.

Task-level productivity studies and studies of human--AI collaboration address different parts of this problem. Experiments on professional writing and evidence from customer support document productivity effects in particular settings \citep{noy,brynjolfsson}; a meta-analysis finds substantial variation in the performance of human--AI combinations relative to their constituent members \citep{vaccaro}. Organization design supplies the intervening objects: the division and allocation of work, dependencies, and integration of effort \citep{puranam,malone}. Complementary organizational investments can also affect the returns to a new technology \citep{milgrom,jcurve}. Consequently, comparing agent models while holding the organization fixed answers a different question from comparing organizations that can revise their workflows.

Workflow revision creates a measurement problem of its own. A new review policy changes both which cases receive attention and which failures become visible. If the same selected cases determine whether that policy is retained, organizational learning can favor a locally successful but globally harmful arrangement. An improvement procedure therefore needs an explicit account of the evidence it can acquire, the changes it can authorize, and the costs it incurs. The procedure may itself become a target of revision.

We call this problem \emph{recursive organization improvement} (ROI; distinct from return on investment): evaluating and retaining changes to organizational arrangements, including changes to the procedures through which those arrangements are improved. The research object is the human--agent organization. Agent capability, human judgment, coordination, evidence access, and authority can change at different rates and under different constraints.

This paper contributes a modeling specification that connects organizational state, actor-visible event histories, and evidence-carrying change contracts. Its distinctive organizing choice is to make the object of a revision and the procedure used to evaluate it jointly explicit. The specification supports reconstructing a decision from its available evidence, checking whether a patch meets declared admission conditions, and comparing retained changes under a stable outcome criterion. These are concrete modeling obligations against which an implementation can be inspected.

We study one mechanism in depth: how an organization discovers, retains, and updates evaluation evidence. The experiment crosses evidence-reset, cumulative, and finite-window memory with fixed evaluation programs, one-time discovery followed by budget reallocation, and repeated program assessment. A template-ranking reversal makes old evidence substantively obsolete. Matching resource ceilings and the acquisition process allows us to distinguish the effect of remembering evidence from the effect of replacing an evaluator. The results show that memory can remove an apparent cost of repeated improvement in a stable environment, yet delay adaptation after change. Additional acquisition-matched controls substantially narrow the gain attributable to program replacement. The specification's role is to make these organizational interventions and their evidence requirements explicit.

The argument follows three stages. Discovery produces information and candidate arrangements. Retention turns observations into reusable organizational knowledge. Reassessment tests whether that knowledge and the arrangements built on it remain useful. These stages connect a stronger agent's task performance to a team's ability to improve its work over time.

\section{Conceptual foundations and modeling requirements}
\label{sec:foundations}
\subsection{Organizational learning and explicit organizations}
Organizational learning connects experience to retained routines \citep{levitt,march}. The distinction between the abstract understanding of a routine and its situated performance helps explain how routines can generate both stability and change \citep{feldman}. Deliberate learning adds articulation and codification of experience \citep{zollo}, while double-loop learning addresses governing assumptions and values \citep{argyris}. Engelbart's improvement-of-improvement distinction supplies a close conceptual antecedent for examining the procedures that generate organizational change \citep{engelbart}. ROI translates this question into versions, events, change targets, and outcome comparisons.

Human--agent organizations also require an account of interdependence. Mixed-initiative interaction and levels of automation distinguish who initiates or performs a function \citep{mixed,levels}; Coactive Design places the support for interdependent activity at the center of system design \citep{coactive}. Appropriate reliance concerns when a person should use automated advice \citep{lee}. These choices alter information exposure, decision rights, and resource use, so they belong in the organizational model rather than being inferred from an actor's capability score.

Existing formalisms already supply much of the representational machinery. MOISE+ models structural, functional, and deontic organization; OperA represents organizational interaction and agent autonomy \citep{moise,opera}. Process mining uses event data for process discovery, conformance checking, and enhancement \citep{processmining}. The proposed specification connects these objects to the evaluation and retention of changes through the relations in \cref{tab:crosswalk}. A sufficiently annotated event log or state-transition model can encode the same relations. The contribution is their explicit organization into a comparison contract, whose completeness and execution can be checked.

\begin{table}[t]
\caption{Connections to established approaches. The right column identifies the records required for an ROI comparison, rather than an expressiveness ranking.}
\label{tab:crosswalk}\small
\begin{tabularx}{\linewidth}{@{}p{.23\linewidth}YY@{}}\toprule
Approach & Established modeling object & Required connection for a change comparison \\\midrule
MOISE+; OperA & Roles, missions, norms, interactions & Target-specific authority, prior version, evidence and comparator for a revision. \\
Process mining & Cases, activities, event order, conformance & Organization/procedure versions, actual observation access, and sampling coverage. \\
Organizational learning & Experience, routines, retention & Link a retained routine to the experiment and criterion used to accept it. \\
Agent/workflow search & Executable candidates and evaluations & Human resources, rights, generated-output costs, and revision of the evaluator. \\
Generic event log & Arbitrary recorded attributes & All the same relations, if explicitly supplied; missing fields remain unknown. \\\bottomrule
\end{tabularx}
\end{table}

Automated agent-system design and workflow optimization provide candidate-generation mechanisms \citep{adas,aflow}; linguistic feedback can support behavioral revision \citep{reflexion}. In ROI, these mechanisms are embedded in an organization with human roles, restricted evidence, and retention decisions. The central comparison concerns organizational outcomes under specified changes, whether the candidate was written by a person, generated by an agent, or selected from a finite catalog.

\subsection{From development practices to measurable arrangements}
Waterfall-like staging, agile development, and Extreme Programming suggest different settings for dependencies, batch size, feedback delay, test placement, and ownership. Royce's original account includes iteration and the risks of late system testing \citep{royce}; agile principles emphasize frequent delivery and reflection \citep{agile}; XP provides more concrete testing, pairing, and integration practices \citep{xp}. These labels identify families of configurations rather than an ordering of organizational quality. A team can retain release gates while shortening internal feedback cycles, or introduce agents into an unchanged approval pipeline.

For example, let drafting, review, and integration capacities be 4, 5, and 9 accepted-task equivalents per day. A deterministic serial capacity bound is their minimum. Doubling drafting alone raises the bound from 4 to 5; raising review capacity to 8 at the same time raises it to 8. The factorial interaction is $8-5-4+4=3$. This standard bottleneck calculation \citep{amdahl,little} isolates why technical and workflow changes can be complementary. Rework, variable quality, and demand must be modeled separately in an empirical application. Our specification records the fields needed to distinguish such joint changes from an increase in task-level speed.
